% pasguide.tex
% v1.0, released 24 Mar 2021
% Copyright 2021 Cambridge University Press

\documentclass{pas}


\usepackage{multirow}

\begin{document}

\lefttitle{Publications of the Astronomical Society of Australia}
\righttitle{Cambridge Author}

\jnlPage{1}{4}
\jnlDoiYr{2021}
\doival{10.1017/pasa.xxxx.xx}

\articletitt{Research Paper}

\title{A Great and Engaging Title about RXJ1720+2638}

\author{\sn{Wellington} \gn{Author1}$^{1}$ and \sn{Wellington} \gn{Author2}$^{2}$}

\affil{$^1$School of Chemical and Physical Sciences, Victoria University of Wellington, PO Box 600, Wellington 6140, New Zealand}

\corresp{A. Cambridge, Email: ACambridge@student.unimelb.edu.au}

\citeauth{Author1 C and Author2 C, an open-source python tool for simulations of source recovery and completeness in galaxy surveys. {\it Publications of the Astronomical Society of Australia} {\bf 00}, 1--12. https://doi.org/10.1017/pasa.xxxx.xx}

\history{(Received xx xx xxxx; revised xx xx xxxx; accepted xx xx xxxx)}

\begin{abstract}
An abstract written to inspire interest and unbroken hope for the future.
\end{abstract}

\begin{keywords}
Key1, Key2, Key3, Key4
\end{keywords}

\maketitle

\section{Introduction}

\begin{itemize}
  \item Introduction to the galaxy cluster RXJ1720+2638 and its general properties, ala Yvette's paper
  \item Introduction to the minihalo, what is a minihalo.
  \item What are the proposed methods to cause said emission, with references to the simulations.
  \item What is this paper focussed on: 8-12 GHz observation of the minihalo. What parts of the paper correspond to what thing.
  \item Mention G14 so that we can mention non-weirdly? during analysis.
  \item Coordinates + cosmological settings assumed etc
\end{itemize}

\section{Observations}

\begin{itemize}
  \item The primary 8-12 GHz dataset info (VLA, configuration, dates, observation proposals, etc)
  \item What primary and secondary calibrators were used.
  \item The L-band observation
\end{itemize}

\section{Data reduction}

\begin{itemize}
  \item Description of the data reduction steps.
  \item Using CASA pipeline to calibrate the data, with a modified flagging step using AOFLAGGER instead of the CASA version.
  \item Mention broadly the steps taken in the pipeline.
  \item Reasoning for above is better quality flagging for the microwave emission in bands, due to Elon Musk or something.
  \item Note where interference takes place, alongside how much data was removed from each observation. 
  \item Imaging done via tclean with broad nterms / weighting parameter. Leave the subtraction and stuff for analysis. 
  \item Separate items for L-band and X-band
\end{itemize}

\section{Analysis}

\subsection{Subtraction of AGN}
\begin{itemize}
  \item Checking the difference in measured flux for the point sources since measurements C \& D are a year apart.
  \item Some statistics on variablity
  \item Uvrange for cutting out diffuse emission + point source modelling.
  \item Show results, compare results of secondary calibrator. 
  \item Conclude its probably fine.
  \item Since variablity exists, each dataset AGN was modelled independently then subtracted.
\end{itemize}
\textbf{Image of point sources, labelled AGN, head-tail, right (better name)}

\subsection{Subtracting the right galaxy thingy}
\begin{itemize}
  \item Difficulty in getting a clean image due to a noisy point source to the right of the minihalo $\sim$ beam size
  \item Also subtracted this using the same process above.
  \item Little information about the source, just mention brightness, size, location.
\end{itemize}

\subsection{Calculating minihalo flux}

\begin{itemize}
  \item With clean image can now calculate the minihalo flux.
  \item In order for comparable results, estimated mask from G14 to use for analysis.
  \item Computed total flux within the regions, then the error. 
  \item Show derivation of the error, justify differences in error calc to G14.
  \item Might as well show stability to robust settings [1, 2] if we have space.
  \item Show possible re-energisation, but below 3x sigma. 
\end{itemize}
\textbf{Image of minihalo with mask contours overlaid}

\subsection{Comparison to other frequencies}
\begin{itemize}
  \item Compare the points above to the G14 data collected + things from Y21.
  \item Show agreement in the full, centre, tail, AGN. Plus comparisons between them.
  \item Lack of steepening.
\end{itemize}
\textbf{Recreation of plot from G14, with additional information}

\subsection{SZ effects}
\begin{itemize}
  \item Look the contribution is negligible.
  \item This is how we worked it out: Little thing is littler than the little thing we are measuring.
\end{itemize}

\subsection{L-band data}
\begin{itemize}
  \item Not entirely sure what to do here.
  \item Also don't have a perfect image yet.
  \item Point at the possible gentle re-energisation also present here.
\end{itemize}
\textbf{Image of L-band minihalo}

\section{Discussion}
\begin{itemize}
  \item Discuss in detail the difference between the hadronic and re-acceleration scenarios.
  \item Talk about the conclusions taken from G14, plus the new data collected since then, in particular the presence of sub-kpc jets Y22 and Y21 results.
  \item More discussion about differences between the tail and centre. 
  \item Actually do some reading.
\end{itemize}

\section{Conclusion}
I like the bullet point style of the main points of the paper, so might as well do that.
\begin{itemize}
  \item Clean imaging of 8-12 GHz VLA observation for RXJ1720+2638 minihalo to measure minihalo flux.
  \item Calculated flux density of the minihalo, tail, centre.
  \item More things from the discussion
\end{itemize}

\begin{thebibliography}{}
\end{thebibliography}

\end{document}

